\chapter{Solution Architecture}
\label{ch:architecture}

generic concepts and open industry standards instead of proprietary concepts/technology unless necessary

Needs to support both OLTP/OLAP use cases on data serving tier 

%%%%%%%%%%%
%%% 3.1 %%%
%%%%%%%%%%%
\section{Architecture Context}

It will be used to pull data from data sources (those are naturally out of scope)

It needs to be ready to support an application security management dashboard with a combination of OLTP and OLAP use cases (dashboard application also out of scope).

%%%%%%%%%%%
%%% 3.2 %%%
%%%%%%%%%%%
\section{Architecture Principles and Constraints}


%%%%%%%%%%%
%%% 3.3 %%%
%%%%%%%%%%%
\section{Technology Selection}

What I'm choosing and why, what were the alternatives and reasons not to choose them

\subsection{Data Engineering Platform}

Choosing \textbf{Databricks} over Snowflake, Fabric, BigQuery

Contains a strong data governance framework - \textbf{Unity Catalog}, good for security non-functional requirement

\subsection{Pipeline Framework}

I will use \gls{dltables}, a "data engineering pipeline framework running on top of Delta Lake that combines incremental ingestion, streamlined ETL, and automated data quality processes such as expectations" \cite{Lee2024}.

The advantages of \gls{dltables} are...

\subsection{Connectors}

(general plan, more details in Integration Architecture)

dlt, Fivetran?, Databricks Partners?

\gls{dltool} is utilized to ingest data from security tools into the Bronze layer. Source code~\ref{code:github-connector} shows an example implementation of a connector for GitHub security data.

\begin{code}{Python}{GitHub Security Connector Implementation}{code:github-connector}
import dlt
from dlt.sources.helpers.rest_client import paginate

@dlt.source
def github_security_source(org: str, token: str):
    """Extract security data from GitHub organization."""

    @dlt.resource(write_disposition="merge", primary_key="id")
    def repos():
        """Repository inventory"""
        yield from paginate(f"/orgs/{org}/repos")

    @dlt.resource(write_disposition="merge", primary_key="number")
    def code_scanning_alerts():
        """GitHub Code Scanning (SAST) alerts"""
        yield from paginate(f"/orgs/{org}/code-scanning/alerts")

    @dlt.resource(write_disposition="merge", primary_key="number")
    def dependabot_alerts():
        """Dependabot (SCA) vulnerability alerts"""
        yield from paginate(f"/orgs/{org}/dependabot/alerts")

    return repos, code_scanning_alerts, dependabot_alerts
\end{code}

Unity Catalog

Delta Lake

Python vs Scala vs SQL


%%%%%%%%%%%
%%% 3.4 %%%
%%%%%%%%%%%
\section{High-Level Solution Architecture}


%%%%%%%%%%%
%%% 3.5 %%%
%%%%%%%%%%%
\section{Component Design}

See Figure~\ref{fig:component-design}

\begin{figure}[htbp]
\centering
\begin{tikzpicture}[
    node distance=0.5cm,
    tier/.style={rectangle, draw=black!40, fill=gray!8, minimum width=11cm, minimum height=1.2cm, rounded corners=3pt},
    tierbox/.style={rectangle, draw=black!60, fill=white, minimum width=2.2cm, minimum height=0.7cm, font=\small, inner sep=5pt, rounded corners=3pt},
    sourcebox/.style={rectangle, draw=black!50, fill=white, minimum width=1.6cm, minimum height=0.6cm, font=\footnotesize, rounded corners=2pt},
    extbox/.style={rectangle, draw=black!50, fill=white, minimum width=1.8cm, minimum height=0.6cm, font=\footnotesize, rounded corners=2pt},
    consumer/.style={rectangle, draw=black!50, fill=gray!15, minimum width=1.6cm, minimum height=0.6cm, font=\footnotesize, rounded corners=2pt},
    arrow/.style={-{Stealth[length=2mm]}, thick, black!60},
    tierlabel/.style={font=\small\bfseries, text=black!70},
    bronze/.style={fill={rgb,255:red,159;green,122;blue,52}, fill opacity=0.35, text=black, text opacity=1},
    silver/.style={fill={rgb,255:red,192;green,192;blue,192}, fill opacity=0.4, text=black, text opacity=1},
    gold/.style={fill={rgb,255:red,212;green,175;blue,55}, fill opacity=0.35, text=black, text opacity=1},
]

% === DATA SOURCES TIER ===
\node[tier] (sources-bg) {};
\node[tierlabel, left, align=right] at (sources-bg.west) {Data\\Sources};
\node[sourcebox] at ($(sources-bg.center) + (-3.6cm,0)$) (github) {GitHub};
\node[sourcebox] at ($(sources-bg.center) + (-1.8cm,0)$) (sonar) {SonarQube};
\node[sourcebox] at ($(sources-bg.center) + (0,0)$) (snyk) {Snyk};
\node[sourcebox] at ($(sources-bg.center) + (1.8cm,0)$) (nvd) {NVD};
\node[sourcebox] at ($(sources-bg.center) + (3.6cm,0)$) (other) {...};

% === EXTERNAL INTEGRATION TOOLS TIER ===
\node[tier, below=0.5cm of sources-bg, fill=blue!5] (external-bg) {};
\node[tierlabel, left, align=right] at (external-bg.west) {Integration\\Tools};
\node[extbox] at ($(external-bg.center) + (-3cm,0)$) (fivetran) {Fivetran};
\node[extbox] at ($(external-bg.center) + (-1cm,0)$) (airbyte) {Airbyte};
\node[extbox] at ($(external-bg.center) + (1cm,0)$) (meltano) {Meltano};
\node[extbox] at ($(external-bg.center) + (3cm,0)$) (custom) {...};

% === DATABRICKS PLATFORM ===
% Ingestion Tier
\node[tier, below=1.1cm of external-bg, minimum height=1.5cm] (ingestion-bg) {};
\node[tierlabel, left, align=right] at (ingestion-bg.west) {Ingestion\\Tier};
\node[tierbox, align=center] at ($(ingestion-bg.center) + (-3.5cm,0)$) (partner-conn) {Partner\\Connectors};
\node[tierbox, align=center] at ($(ingestion-bg.center) + (0,0)$) (dlt-conn) {dlt\\Connectors};
\node[tierbox, align=center] at ($(ingestion-bg.center) + (3.5cm,0)$) (custom-api) {Custom\\APIs};

% Processing Tier
\node[tier, below=0.5cm of ingestion-bg, minimum height=1.5cm] (processing-bg) {};
\node[tierlabel, left, align=right] at (processing-bg.west) {Processing\\Tier};
\node[tierbox, bronze, minimum width=2.8cm] at ($(processing-bg.center) + (-3.5cm,0)$) (bronze) {Bronze};
\node[tierbox, silver, minimum width=2.8cm] at (processing-bg.center) (silver) {Silver};
\node[tierbox, gold, minimum width=2.8cm] at ($(processing-bg.center) + (3.5cm,0)$) (gold) {Gold};

% Serving Tier
\node[tier, below=0.5cm of processing-bg, minimum height=1.5cm] (serving-bg) {};
\node[tierlabel, left, align=right] at (serving-bg.west) {Serving\\Tier};
\node[tierbox] at ($(serving-bg.center) + (-1.8cm,0)$) (lakebase) {Lakebase};
\node[tierbox] at ($(serving-bg.center) + (1.8cm,0)$) (restapi) {REST API};

% === DATA CONSUMERS ===
\node[tier, below=0.9cm of serving-bg, fill=gray!15] (consumers-bg) {};
\node[tierlabel, left] at (consumers-bg.west) {Data Consumers};
\node[consumer, fill=white] at ($(consumers-bg.center) + (-2.5cm,0)$) (jira) {Jira};
\node[consumer, fill=white] at ($(consumers-bg.center) + (0,0)$) (aspm) {ASPM};
\node[consumer, fill=white] at ($(consumers-bg.center) + (2.5cm,0)$) (siem) {SIEM};

% === ARROWS ===
% Sources to External Tools
\draw[arrow] (sources-bg.south) -- (external-bg.north);

% External Tools to Ingestion
\draw[arrow] (external-bg.south) -- (ingestion-bg.north);

% Direct bypass arrow (Sources to Ingestion) - goes outside on the right
\draw[arrow] (sources-bg.east) -- ++(0.8cm,0) -- ($(ingestion-bg.east) + (0.8cm,0)$) -- (ingestion-bg.east);

% Ingestion to Processing
\draw[arrow] (ingestion-bg.south) -- (processing-bg.north);

% Processing flow
\draw[arrow] (bronze.east) -- (silver.west);
\draw[arrow] (silver.east) -- (gold.west);

% Processing to Serving
\draw[arrow] (processing-bg.south) -- (serving-bg.north);

% Serving to Consumers
\draw[arrow] (serving-bg.south) -- (consumers-bg.north);

% === DATABRICKS BOUNDARY ===
\begin{scope}[on background layer]
\node[draw=red!30, fill=red!5, rounded corners=5pt, fit=(ingestion-bg)(processing-bg)(serving-bg), inner sep=8pt, label={[font=\footnotesize, text=black]above:Databricks Platform}] {};
\end{scope}

\end{tikzpicture}
\caption{Component Design}
\label{fig:component-design}
\end{figure}


\subsection{Data Ingestion Tier}

The various connectors

\subsection{Data Processing Tier}

\subsubsection{Pipeline Architecture}

Medallion architecture, see Chapter~\ref{ch:design}

\gls{dltables} for defining transformations

\subsubsection{Pipeline Orchestration}

Databricks Workflows, don't think my workflows will be complex enough to use Airflow DAGs

\subsection{Data Serving Tier}

Lakebase

%%%%%%%%%%%
%%% 3.6 %%%
%%%%%%%%%%%
\section{Integration Architecture}


\subsection{Connector Strategy}

Different patterns and tools optimal for different data sources

As a service: Fivetran if it satisfies requirements (easiest to do, no need to reinvent the wheel)

Reuse existing: e.g. from DefectDojo https://github.com/DefectDojo/django-DefectDojo/tree/master/dojo/tools

Databricks Partner

dlt as default tool for custom connectors to REST APIs

\subsubsection{Application Inventory}

\subsubsection{ServiceNow CMDB}

Call ServiceNow API vs build ServiceNow app



\subsubsection{}

%%%%%%%%%%%
%%% 3.7 %%%
%%%%%%%%%%%
\section{Extensibility Architecture}

Add new integrations via API

%%%%%%%%%%%
%%% 3.8 %%%
%%%%%%%%%%%
\section{Security Architecture}